\documentclass[12pt, letterpaper]{article}

\usepackage[utf8]{inputenc}
\usepackage[T1]{fontenc}
\usepackage[spanish]{babel}
\usepackage{geometry}
\usepackage{amsmath, amssymb, amsthm}
\usepackage{booktabs}
\usepackage{array}
\usepackage{microtype}
\usepackage{setspace}
\usepackage{parskip}
\usepackage{titlesec}
\usepackage{fancyhdr}
\usepackage{enumitem}
\usepackage{bm}
\usepackage{mathtools}
\usepackage{tcolorbox}
\usepackage{hyperref}

\tcbuselibrary{skins, breakable}

\geometry{top=2.8cm, bottom=3.0cm, left=3.2cm, right=3.2cm}
\setstretch{1.20}

\pagestyle{fancy}
\fancyhf{}
\fancyhead[C]{\small\textsc{Econometria · Gujarati --- Ejercicio 10.6}}
\fancyfoot[C]{\small\thepage}
\renewcommand{\headrulewidth}{0.4pt}
\renewcommand{\footrulewidth}{0pt}

\titleformat{\section}
  {\large\bfseries\scshape}{\thesection.}{0.6em}{}
  [\vspace{-4pt}\rule{\linewidth}{0.4pt}]
\titleformat{\subsection}{\normalsize\bfseries}{\thesubsection.}{0.5em}{}
\titleformat{\subsubsection}{\normalsize\itshape}{}{0em}{}

\newtheorem{prop}{Proposicion}
\newtheorem{lema}{Lema}
\theoremstyle{remark}
\newtheorem{obs}{Observacion}

\newtcolorbox{clave}{
  colback=white, colframe=black,
  boxrule=0.5pt, arc=3pt,
  left=8pt, right=8pt, top=6pt, bottom=6pt,
  breakable
}

\newcommand{\E}{\mathbb{E}}
\newcommand{\Var}{\operatorname{Var}}
\newcommand{\Cov}{\operatorname{Cov}}
\newcommand{\plim}{\operatorname{plim}}
\newcommand{\bhat}{\hat{\beta}}

\begin{document}

\begin{center}
  {\LARGE\bfseries Propension Marginal a Consumir,}\\[6pt]
  {\LARGE\bfseries Multicolinealidad y Sesgo por Variable Omitida}\\[12pt]
  {\large\itshape Interpretacion de la diferencia entre los modelos (10.6.1) y (10.6.4)}\\[14pt]
  \rule{0.55\linewidth}{0.6pt}\\[8pt]
  {\normalsize Ejercicio~10.6 --- \textit{Econometria}, Gujarati \& Porter, 5.a ed.}\\[4pt]
  {\small\textsc{Derivacion completa del sesgo, algebra matricial y diagnostico de colinealidad}}
\end{center}

\vspace{10pt}

\noindent\textbf{Contexto.}\quad
El ejemplo~10.1 de Gujarati estima una funcion de consumo donde la variable
dependiente es el gasto de consumo personal $Y_i$ y los regresores son el
ingreso personal disponible $X_{2i}$ y la riqueza personal $X_{3i}$. Se
presentan dos versiones del modelo: una con ambas variables y otra eliminando
la riqueza. La propension marginal a consumir (PMC) estimada difiere
sustancialmente entre los dos modelos, lo que plantea una pregunta de
interpretacion econometrica fundamental.

\section{Los dos modelos en cuestion}

\subsection{Modelo completo: ingreso y riqueza (10.6.1)}

\begin{equation}
  \hat{Y}_i = \hat{\beta}_1 + \hat{\beta}_2 X_{2i} + \hat{\beta}_3 X_{3i}
  \label{eq:completo}
\end{equation}

\noindent con los siguientes resultados estimados:

\begin{table}[ht]
\centering
\caption{Resultados del modelo completo~(\ref{eq:completo})}
\label{tab:completo}
\small
\begin{tabular}{lrrrr}
\toprule
\textbf{Variable} & \textbf{Coeficiente} & \textbf{Error estandar}
                  & \textbf{Estadistico $t$} & \textbf{$p$-valor} \\
\midrule
Constante         & $-184.0$  & ---      & ---     & ---    \\
Ingreso $X_2$     & $0.9415$  & $0.8229$ & $1.14$  & $> 0.05$ \\
Riqueza $X_3$     & $-0.0424$ & $0.0807$ & $-0.52$ & $> 0.05$ \\
\midrule
\multicolumn{2}{l}{$R^2 = 0.9608$} & \multicolumn{3}{l}{$F$-estadistico: significativo} \\
\bottomrule
\end{tabular}
\end{table}

\noindent\textbf{Observacion critica}: ninguno de los coeficientes individuales
es estadisticamente significativo al nivel convencional del $5\%$, a pesar de
que el $R^2$ global es $0.9608$. Esto es la firma clasica de multicolinealidad
severa.

\subsection{Modelo reducido: solo ingreso (10.6.4)}

\begin{equation}
  \hat{Y}_i = \hat{\beta}_1^* + \hat{\beta}_2^* X_{2i}
  \label{eq:reducido}
\end{equation}

\noindent con los siguientes resultados:

\begin{table}[ht]
\centering
\caption{Resultados del modelo reducido~(\ref{eq:reducido})}
\label{tab:reducido}
\small
\begin{tabular}{lrrrr}
\toprule
\textbf{Variable} & \textbf{Coeficiente} & \textbf{Error estandar}
                  & \textbf{Estadistico $t$} & \textbf{$p$-valor} \\
\midrule
Constante         & $-184.0$  & ---      & ---     & ---    \\
Ingreso $X_2$     & $0.5091$  & $0.0357$ & $14.24$ & $< 0.001$ \\
\midrule
\multicolumn{2}{l}{$R^2 = 0.9621$} & \multicolumn{3}{l}{} \\
\bottomrule
\end{tabular}
\end{table}

\noindent Al eliminar $X_3$, el coeficiente de $X_2$ cae de $0.9415$ a
$0.5091$ y su estadistico $t$ salta de $1.14$ a $14.24$.

\section{La multicolinealidad como causa de la discrepancia}

\subsection{Diagnostico de colinealidad}

La regresion auxiliar de $X_3$ sobre $X_2$ produce un $R^2$ de $0.9979$.
El factor de inflacion de varianza (VIF) del ingreso en el modelo completo es:

\begin{equation}
  \text{VIF}(X_2) = \frac{1}{1 - R_{2\cdot 3}^2}
                  = \frac{1}{1 - 0.9979}
                  = \frac{1}{0.0021}
                  \approx 476
  \label{eq:VIF}
\end{equation}

\noindent Un VIF de $476$ es extremo: la varianza del estimador de $\beta_2$
es $476$ veces mayor de lo que seria si $X_2$ y $X_3$ fuesen ortogonales.
Esto explica por que el error estandar de $\hat\beta_2$ en~(\ref{eq:completo})
es $0.8229$ ---enorme--- mientras que en~(\ref{eq:reducido}) es solo $0.0357$.

\subsection{Efecto sobre los errores estandar}

La varianza del estimador MCO de $\beta_2$ en el modelo completo es:

\begin{equation}
  \Var(\hat\beta_2) = \frac{\sigma^2}{\sum x_{2i}^2} \cdot \frac{1}{1 - r_{23}^2}
  \label{eq:var_beta2}
\end{equation}

\noindent donde $x_{2i} = X_{2i} - \bar X_2$ y $r_{23}$ es la correlacion entre
$X_2$ y $X_3$. Con $r_{23}^2 = 0.9979$, el factor $1/(1 - r_{23}^2) = 476$
infla la varianza en ese factor respecto al modelo sin $X_3$. El error estandar
escala con la raiz de este factor:

\begin{equation}
  \frac{\text{ee}(\hat\beta_2)^{\text{completo}}}{\text{ee}(\hat\beta_2)^{\text{reducido}}}
  \approx \sqrt{476} \approx 21.8
  \label{eq:razon_ee}
\end{equation}

\noindent Efectivamente, $0.8229 / 0.0357 \approx 23$, coherente con este calculo.

\section{El sesgo por variable omitida: derivacion formal}

\subsection{Planteamiento del problema}

Suponga que el verdadero modelo de la economia es el modelo completo:

\begin{equation}
  Y_i = \beta_1 + \beta_2 X_{2i} + \beta_3 X_{3i} + u_i
  \label{eq:verdadero}
\end{equation}

\noindent Pero el econometrista estima en cambio el modelo reducido~(\ref{eq:reducido}),
omitiendo $X_3$. Sea $b_{12}$ el estimador de la pendiente de la regresion de $Y$
sobre $X_2$ unicamente, y $b_{32}$ el estimador de la pendiente de la regresion de
$X_3$ sobre $X_2$:

\begin{align}
  b_{12} &= \frac{\sum x_{2i} y_i}{\sum x_{2i}^2}
  \label{eq:b12}\\[4pt]
  b_{32} &= \frac{\sum x_{2i} x_{3i}}{\sum x_{2i}^2}
  \label{eq:b32}
\end{align}

\noindent donde las minusculas indican desviaciones respecto a la media.

\subsection{Demostracion del sesgo}

Sustituimos~( \ref{eq:verdadero} ) en~(\ref{eq:b12}):

\begin{align}
  b_{12}
  &= \frac{\sum x_{2i}(\beta_2 x_{2i} + \beta_3 x_{3i} + u_i)}{\sum x_{2i}^2}
  \notag\\[4pt]
  &= \beta_2 \cdot \frac{\sum x_{2i}^2}{\sum x_{2i}^2}
   + \beta_3 \cdot \frac{\sum x_{2i} x_{3i}}{\sum x_{2i}^2}
   + \frac{\sum x_{2i} u_i}{\sum x_{2i}^2}
  \notag\\[4pt]
  &= \beta_2 + \beta_3 b_{32} + \frac{\sum x_{2i} u_i}{\sum x_{2i}^2}
  \label{eq:sesgo_derivacion}
\end{align}

Tomando esperanza (bajo el supuesto de que $X_2$ es exogena, $\E[\sum x_{2i} u_i] = 0$):

\begin{equation}
  \boxed{
    \E[b_{12}] = \beta_2 + \beta_3 \cdot b_{32}
  }
  \label{eq:sesgo}
\end{equation}

\noindent El termino $\beta_3 \cdot b_{32}$ es el \textbf{sesgo por variable omitida}.
El estimador $b_{12}$ es sesgado e inconsistente siempre que $\beta_3 \neq 0$
y $b_{32} \neq 0$.

\subsection{Calculo numerico del sesgo}

Del ejercicio tenemos los siguientes valores:

\begin{align}
  \hat\beta_2 &= 0.9415 \quad \text{(coeficiente de ingreso en modelo completo)}
  \label{eq:valores1}\\
  \hat\beta_3 &= -0.0424 \quad \text{(coeficiente de riqueza en modelo completo)}
  \label{eq:valores2}\\
  b_{32}      &= 10.191 \quad \text{(pendiente de } X_3 \text{ regresado sobre } X_2\text{)}
  \label{eq:valores3}\\
  b_{12}      &= 0.5091 \quad \text{(PMC del modelo reducido)}
  \label{eq:valores4}
\end{align}

El sesgo es:

\begin{equation}
  \text{Sesgo} = \hat\beta_3 \cdot b_{32} = (-0.0424)(10.191) = -0.4321
  \label{eq:sesgo_numerico}
\end{equation}

\noindent Verificamos la identidad~(\ref{eq:sesgo}) con los datos:

\begin{equation}
  b_{12} = \hat\beta_2 + \hat\beta_3 \cdot b_{32}
         = 0.9415 + (-0.4321)
         = 0.5094
         \approx 0.5091 \quad \checkmark
  \label{eq:verificacion}
\end{equation}

\noindent La pequena discrepancia se debe al redondeo. La identidad se cumple
exactamente en algebra de precision infinita.

\begin{clave}
  \textbf{Interpretacion.} El $0.5091$ del modelo reducido no es la PMC
  ``verdadera''. Es una mezcla del efecto directo del ingreso ($\beta_2$) y
  del efecto indirecto de la riqueza que el ingreso captura en su ausencia
  ($\beta_3 b_{32}$). Dado que $b_{32} = 10.191 > 0$ (riqueza y ingreso se
  mueven juntos) y $\beta_3 = -0.0424 < 0$ (riqueza tiene efecto negativo),
  el sesgo es negativo: omitir la riqueza subestima la PMC verdadera.
\end{clave}

\section{Interpretacion economica de la diferencia}

\subsection{Por que difieren tanto los coeficientes}

La diferencia $0.9415 - 0.5091 = 0.4324 \approx |\text{sesgo}|$ se explica
completamente por la formula~(\ref{eq:sesgo}). No hay misterio econometrico:
es una consecuencia algebraica de omitir una variable correlacionada con el
regresor incluido.

\subsection{Que mide cada coeficiente}

\begin{enumerate}[leftmargin=1.8em, itemsep=6pt]
  \item \textbf{$\hat\beta_2 = 0.9415$ en el modelo completo} es el
    \textbf{efecto parcial} del ingreso sobre el consumo, manteniendo constante
    la riqueza. Mide cuanto cambia $Y$ ante un cambio de una unidad en $X_2$
    \emph{sin mover} $X_3$. El problema es que con $r_{23}^2 = 0.9979$, los
    datos no tienen variacion suficiente de $X_2$ con $X_3$ constante:
    el error estandar es tan grande ($0.8229$) que el intervalo de confianza
    al $95\%$ para $\beta_2$ abarca desde valores negativos hasta casi~$2.6$.
    Es decir, el coeficiente es estadisticamente inutil, no porque sea economicamente
    cero, sino porque la colinealidad impide estimarlo.

  \item \textbf{$b_{12} = 0.5091$ en el modelo reducido} es el
    \textbf{efecto total} del ingreso sobre el consumo, incluyendo su correlacion
    con la riqueza. Cuando el ingreso sube en una unidad, la riqueza tambien
    tiende a subir en $b_{32} = 10.191$ unidades; ese movimiento de la riqueza
    tambien afecta al consumo con coeficiente $\beta_3 = -0.0424$. El modelo
    reducido atribuye ese efecto conjunto al ingreso, por eso su coeficiente
    es diferente.
\end{enumerate}

\subsection{Signo incorrecto del coeficiente de riqueza}

En el modelo completo, $\hat\beta_3 = -0.0424$: la riqueza tiene un efecto
negativo sobre el consumo. Esto contradice la teoria economica (mayor riqueza
deberia inducir mayor consumo). Este signo incorrecto es tambien consecuencia
de la colinealidad: con VIF $= 476$, los estimadores son extremadamente sensibles
a pequeñas perturbaciones en los datos. El verdadero valor de $\beta_3$ muy
probablemente es positivo, pero la colinealidad genera una estimacion con signo
erroneo.

\subsection{Paradoja del alto $R^2$ con bajos $t$}

Esta situacion ---$R^2 = 0.9608$ con $t_{X_2} = 1.14$ y $t_{X_3} = -0.52$---
es el ejemplo canonico de multicolinealidad. Se explica porque:

\begin{equation}
  R^2_{\text{conjunto}} \neq \sum_j \text{(contribucion individual de } X_j\text{)}
  \label{eq:paradoja}
\end{equation}

El ingreso y la riqueza juntos explican el $96\%$ de la variacion del consumo,
pero sus efectos individuales son tan similares que el modelo no puede discernir
cual de los dos es responsable de que fraccion. El test $F$ global es significativo
(rechaza $\beta_2 = \beta_3 = 0$) pero los tests $t$ individuales no pueden
rechazar $\beta_2 = 0$ ni $\beta_3 = 0$ por separado.

\section{Sintesis y conclusion}

La diferencia entre $0.9415$ y $0.5091$ no refleja un comportamiento distinto
del consumidor. Refleja tres fenomenos econometricos interrelacionados:

\begin{enumerate}[leftmargin=1.8em, itemsep=6pt]
  \item \textbf{Multicolinealidad casi perfecta} ($r_{23}^2 = 0.9979$,
    VIF $= 476$): los datos no tienen variacion independiente de ingreso y
    riqueza, haciendo imposible separar sus efectos con precision.

  \item \textbf{Inestabilidad del modelo completo}: el coeficiente $0.9415$
    tiene un error estandar de $0.8229$, lo que lo hace estadisticamente
    indistinguible de cero. La estimacion es tecnicamente valida pero
    practicamente inutilizable.

  \item \textbf{Sesgo por variable omitida en el modelo reducido}: el
    coeficiente $0.5091$ combina el efecto propio del ingreso con parte
    del efecto de la riqueza. El sesgo exacto es $(-0.0424)(10.191) = -0.4321$,
    que sumado a $0.5091$ recupera $0.9412 \approx 0.9415$.
\end{enumerate}

\begin{clave}
  Ninguno de los dos coeficientes es la PMC ``verdadera'' en un sentido
  incondicional. El $0.9415$ es el efecto parcial del ingreso con riqueza
  constante, pero es impreciso. El $0.5091$ es el efecto total del ingreso
  (absorbiendo el canal de riqueza), pero esta sesgado si la riqueza es
  causalmente relevante. La solucion correcta no es elegir entre los dos
  modelos, sino conseguir datos con suficiente variacion independiente de
  ambas variables para estimar el modelo completo con precision.
\end{clave}

\appendix

\section{Demostracion general del sesgo por variable omitida}

Sea el modelo verdadero con $k$ regresores:
\begin{equation}
  \mathbf{y} = \mathbf{X}_1\boldsymbol{\beta}_1 + \mathbf{X}_2\boldsymbol{\beta}_2 + \mathbf{u}
  \label{eq:gral_completo}
\end{equation}

El estimador MCO del modelo mal especificado (omitiendo $\mathbf{X}_2$) es:
\begin{equation}
  \hat{\boldsymbol{\beta}}_1^{\text{red}}
    = (\mathbf{X}_1'\mathbf{X}_1)^{-1}\mathbf{X}_1'\mathbf{y}
  \label{eq:MCO_red}
\end{equation}

Sustituyendo~(\ref{eq:gral_completo}) en~(\ref{eq:MCO_red}):
\begin{align}
  \hat{\boldsymbol{\beta}}_1^{\text{red}}
  &= (\mathbf{X}_1'\mathbf{X}_1)^{-1}\mathbf{X}_1'
     (\mathbf{X}_1\boldsymbol{\beta}_1 + \mathbf{X}_2\boldsymbol{\beta}_2 + \mathbf{u})
  \notag\\
  &= \boldsymbol{\beta}_1
   + \underbrace{(\mathbf{X}_1'\mathbf{X}_1)^{-1}\mathbf{X}_1'\mathbf{X}_2}_{\mathbf{B}_{21}}
     \boldsymbol{\beta}_2
   + (\mathbf{X}_1'\mathbf{X}_1)^{-1}\mathbf{X}_1'\mathbf{u}
  \label{eq:sesgo_gral}
\end{align}

Tomando esperanza:
\begin{equation}
  \E[\hat{\boldsymbol{\beta}}_1^{\text{red}}]
    = \boldsymbol{\beta}_1 + \mathbf{B}_{21}\boldsymbol{\beta}_2
  \label{eq:esperanza_sesgo}
\end{equation}

\noindent donde $\mathbf{B}_{21} = (\mathbf{X}_1'\mathbf{X}_1)^{-1}\mathbf{X}_1'\mathbf{X}_2$
es la matriz de coeficientes de la proyeccion de $\mathbf{X}_2$ sobre $\mathbf{X}_1$.
El sesgo es $\mathbf{B}_{21}\boldsymbol{\beta}_2$.

Para el caso escalar del ejercicio ($\mathbf{X}_1 = X_2$, $\mathbf{X}_2 = X_3$):
\begin{equation}
  B_{21} = b_{32} = \frac{\sum x_{2i} x_{3i}}{\sum x_{2i}^2} = 10.191
\end{equation}
\begin{equation}
  \text{Sesgo} = b_{32}\beta_3 = (10.191)(-0.0424) = -0.4321 \qquad \square
\end{equation}

\section{Condiciones para que el sesgo sea cero}

El estimador $b_{12}$ es insesgado si y solo si se cumple al menos una de:

\begin{enumerate}[leftmargin=1.8em, itemsep=4pt]
  \item $\beta_3 = 0$: la variable omitida no afecta a $Y$ (irrelevancia).
  \item $b_{32} = 0$: la variable omitida es ortogonal al regresor incluido
    ($X_3 \perp X_2$).
\end{enumerate}

En el presente ejercicio no se cumple ninguna: $\beta_3 = -0.0424 \neq 0$ y
$b_{32} = 10.191 \neq 0$, lo que garantiza un sesgo de magnitud $0.4321$.

\vspace{14pt}
\noindent\rule{\linewidth}{0.4pt}

\noindent{\small\textit{Nota.} Los valores numericos provienen de la tabla de datos
del ejemplo~10.1 de Gujarati \& Porter (2010). La formula de sesgo~(\ref{eq:sesgo})
es exacta en muestras finitas bajo exogeneidad de $X_2$; bajo endogeneidad se
requiere la formula de sesgo asintotico usando probabilidades limite.}

\end{document}
